% \subsection{Gradient-Based Machine Unlearning}
% Weight-editing methods erase target facts by modifying model parameters directly, motivated by privacy regulations and the right to erasure~\cite{yao2024large}. Early approaches relied on prompting-based or in-context interventions~\cite{liu2024survey}; Łucki et al.~\cite{lucki2025adversarial} show that such inference-time methods leave parametric knowledge intact, with models remaining vulnerable to adversarial extraction under targeted queries. Parameter-editing methods address this gap but introduce a different challenge: none accounts for variation in how deeply individual facts are encoded across the model.

% Gradient Ascent (GA) maximizes the negative log-likelihood on the forget set, which frequently leads to catastrophic collapse of linguistic coherence~\cite{yao2024large, zhang2024negative}. Gradient Difference (GD) adds a retain loss to bound collateral damage but treats every token in the forget set with equal importance~\cite{bu2025unlearning}. Negative Preference Optimization (NPO) frames unlearning as preference alignment, pushing the output distribution away from the memorised response~\cite{zhang2024negative}. Weighted Gradient Ascent (WGA) weights each token's gradient by the model's log-probability, concentrating unlearning effort on high-confidence tokens~\cite{wang2025rethinking}. All four methods calibrate the unlearning signal from the model's current output distribution, which changes throughout training and does not reflect how deeply a fact is parametrically encoded~\cite{zhang2024negative, wang2025rethinking}. Methods that apply a retain loss (GD, WGA) further require manual selection of the retain coefficient $\alpha$: too small and general utility degrades; too large and target facts survive erasure. AdaPop replaces this local signal with a global popularity proxy and eliminates manual retain tuning with a dual-ascent controller that adjusts $\alpha$ automatically, directly addressing both the per-fact calibration gap and the hyperparameter sensitivity that none of these methods resolves.

% \subsection{Knowledge Popularity and Memorization}
% Kandpal et al.~\cite{kandpal2023large} show that LLMs recall long-tail facts far less reliably than frequent ones, and Hartmann et al.~\cite{hartmann2024undesirable} survey how parametric memorisation scales with training-corpus frequency. Popular facts are encoded redundantly across many layers; rare facts are lightly encoded and susceptible to over-erasure when subjected to the same gradient force. This asymmetry means a single gradient weight cannot address both regimes simultaneously.

% Log-probabilities do not accurately reflect memorisation depth~\cite{zhang2024negative}: a model may assign high confidence to a surface token pattern without the corresponding fact being deeply encoded across layers. External popularity proxies, such as entity citelink counts or co-occurrence frequencies, provide a fact-level signal independent of the model's current output distribution {\cite{krishnannot, anna2026anatomy}}. AdaPop assigns each fact a popularity-dependent exponent, treating memorisation depth as a property of the fact rather than of the model.
